\clearpage
\section{结论}
本文围绕 Global RRP 的年度滚动参数配置，研究多资产风险预算、收益短缺与实施规则之间的联系。风险预算参考由严格凸的对数问题构造，最终权重由跟踪距离、归一化方差及预测收益短缺组成的凸目标求得。收益目标取参考组合的当期预测收益，多头单位预算使组合保持无杠杆。

数据与执行设计明确区分资产池和当期合格集合。历史观察不足的基金不能提前参与，已合格资产允许形成零权重。均值采用指数加权有效观察，协方差使用共同完整历史并估计收缩强度。交易由目标与漂移权重之差确定，日净收益扣除统一约定成本，周度记录分别保存自然周和持有期收益。

在共同评价区间内，Global RRP 净年化收益为 \globalNetReturn，最大回撤为 \globalMaxDD，夏普为 \globalSharpe。与四个对照相比，主模型处于低收益低波动组合和高收益高波动组合之间，未在所有指标上占优。换手与成本仍是理解结果的必要部分。

\primaryRebalances{} 次周度决策均保留核验记录，\primaryAssetsUsed{} 只 ETF 曾在合格期间形成实质持仓。参数验证、日收益、成本和净值均通过独立重算。十个年度中，九组候选每年都进入一倍标准误集合，说明当前样本没有精确识别惩罚系数。这些证据支持实现的可重复性，不支持唯一参数结论或未来收益保证。

后续研究应冻结当前校准程序和评价口径，并使用新增数据继续检验。实际交易成本、市场容量和机构负债需要独立输入。保持规则连续，才能判断当前收益风险关系能否延续。本文将公开结果限定在指定配置，历史研究记录不承担当前模型的绩效证明功能。
\label{sec:body_end}
